% 書末：參考資料與延伸閱讀（不編章號）
\chapter*{參考資料與延伸閱讀}
\addcontentsline{toc}{chapter}{參考資料與延伸閱讀}
\markboth{參考資料與延伸閱讀}{參考資料與延伸閱讀}
\label{ch:references}

各章最後的「延伸閱讀」提到的書，完整的書目資料都列在這裡，每本附一兩句它適合
做什麼，以及和本書哪幾章搭配。 括號裡的章節號碼是那本書的章節，以這裡列出的
版本為準；換了版本，章節號碼可能不同。

本書的講解、例子和程式都是自己寫的，參考這些書的地方是題型的分類、講解的先後
和經典的結論。 想看另一種講法、更嚴謹的證明，或想多做一些題目時，可以從下面
挑一本。

\section*{篇幅短、適合快速複習}

\begin{itemize}
\item Antti Laaksonen，\textit{Competitive Programmer's Handbook}，2018 年 7 月
3 日草稿版，作者公開的 PDF。

每個主題只用幾頁、一兩個小例子講完，適合讀完本書某一章後，花十幾分鐘從另一個
角度再看一次，也適合想快速看過全貌的人。 本書第~1~到第~13~章的延伸閱讀幾乎
都引用了它：第~1~章（第 1、2 章），第~2~章（第 3 章），第~3~章（3.3 節），
第~4~章（第 5 章），第~5~章（2.4、8.1、9.1 節），第~6~章（4.5、8.2、8.3 節），
第~7~章（第 6 章），第~8~章（3.1 節），第~9~章（第 7 章），第~10~章（7.3 節），
第~11~章（第 11、12 章），第~12~章（第 12 章），第~13~章（第 14 章）。
\end{itemize}

\section*{題型分類細、題目多}

\begin{itemize}
\item 吳邦一，《AP325：從 APCS 實作題檢測三級到五級》，v1.5，2024 年 9 月，作者
公開的講義。

以 APCS 實作題的程度為目標，每個主題附很多例題和習題，還有測試資料，適合要準
備 APCS 的讀者多練。 它的第 0 章（預備知識）到第 8 章（樹上演算法）和本書第~1~%
到第~13~章的範圍大致平行：遞迴、排序與二分搜、佇列與堆疊、貪心與掃描線、分治、
動態規劃、圖、樹都有。 各章引用的是第~13~章（第 8 章）和第~14~章（第 0 章）。

\item Steven Halim、Felix Halim、Suhendry Effendy，\textit{Competitive
Programming 4}，Book 1 與 Book 2，2020（作者自行出版）。

把題目依技巧分得很細，每一節後面附一長串分好類的練習題，適合打算參加程式競
賽、想大量練習的讀者。 第 1 章教怎麼從題目的限制判斷題型、估計複雜度，和本書
第~1、14~章的想法相同。 和本書搭配的地方：第~1~章（Book 1 的 1.3 節），
第~3~章（Book 1 的 3.3.1 節），第~5~章（Book 1 的 3.5.2 節），第~10~章（Book 1
的 3.5.2 節），第~12~章（Book 1 的 4.2.5 節，Book 2 的 8.2.2 節），第~14~章
（Book 1 的 1.3 節）。

\item 劉汝佳，《算法競賽入門經典》，第 2 版，清華大學出版社，2014。

簡體中文。 例題多，程式寫得精簡，適合想用中文多看幾種題型的讀者。 和本書搭配
的地方：第~4~章（第 7 章），第~6~章（6.1 節），第~9~章（9.1、9.2 節），第~10~%
章（9.1 節），第~12~章（6.4、7.5 節）。
\end{itemize}

\section*{嚴謹的證明與演算法設計}

\begin{itemize}
\item Thomas H. Cormen、Charles E. Leiserson、Ronald L. Rivest、Clifford Stein，
\textit{Introduction to Algorithms}，第 4 版，The MIT Press，2022。

大學演算法課常用的教科書，每個演算法都有完整的正確性證明和分析。 想知道本書
省略的證明細節，或要修演算法課的讀者，可以對照著讀。 和本書搭配的地方：第~2~章
（2.1、2.3 節，第 8 章），第~5~章（16.1 節），第~6~章（第 6 章，16.1 節），第~7~%
章（15.1、15.2 節），第~8~章（2.3、4.4 節），第~9~章（14.1、14.3 節），第~10~%
章（14.4 節），第~11~章（20.1 到 20.3 節），第~13~章（10.3、12.1 節）。

\item Steven S. Skiena，\textit{The Algorithm Design Manual}，第 3 版，Springer，
2020。

重點放在怎麼想出演算法、怎麼替一條看似合理的規則找反例，以及怎麼把實際的問題
對應到已知的問題，後半是一份常見問題的目錄。 適合想理解「做法是怎麼想出來的」
的大學生和自學的人。 和本書搭配的地方：第~1~章（第 2 章），第~2~章（4.1 節），
第~3~章（5.1 節），第~7~章（1.2、1.3 節），第~8~章（1.4 節，第 5 章），第~11~章
（第 7 章），第~14~章（第 13 章）。
\end{itemize}

\section*{本系列的其他書}

\begin{itemize}
\item 《C++ 程式設計：給學過 C 的人》：各章提到的「銜接書」，本書用到的 C++ 語法
和 STL 都在那裡。
\item 《C++ 演算法解題 2：動態規劃與圖論》：本書的下一本。 背包、子序列、區間、
位元、數位等動態規劃，拓樸排序、並查集、加權最短路徑、最小生成樹和樹形 DP 都
在那裡。
\end{itemize}
